\documentclass[14pt,a4paper]{extarticle}
\usepackage[T2A]{fontenc}
\usepackage[utf8]{inputenc}
\usepackage[english,russian]{babel}
\usepackage{lmodern}
\usepackage{geometry}
\usepackage{array}
\usepackage{tabularx}
\usepackage{longtable}
\usepackage{ltablex}
\usepackage{booktabs}
\usepackage{hyperref}
\usepackage{indentfirst}
\usepackage{microtype}
\usepackage{graphicx}
\usepackage{setspace}
\usepackage{listings}
\usepackage{graphicx}
\usepackage{caption}
\usepackage{float}
\usepackage{tocloft}
\usepackage{titlesec}
\usepackage{newtxtext,newtxmath}

\titleformat{\section}        % форматируем \section
  {\normalfont\bfseries}      % стиль: обычный + жирный
  {\thesection.}              % номер секции с точкой
  {2mm}                       % расстояние между номером и текстом
  {}                           % перед заголовком можно добавить что-то
\titlespacing*{\section}{1.5cm}{*3}{*1}  % {отступ слева}{отступ сверху}{отступ снизу}

\titleformat{\subsection}        % форматируем \section
  {\normalfont\bfseries}      % стиль: обычный + жирный
  {\thesubsection.}              % номер секции с точкой
  {2mm}                       % расстояние между номером и текстом
  {}                           % перед заголовком можно добавить что-то
\titlespacing*{\subsection}{1.5cm}{*3}{*1}  % {отступ слева}{отступ сверху}{отступ снизу}

\captionsetup{justification=centering}
\renewcommand{\ttdefault}{cmtt}
\keepXColumns
\setlength{\parindent}{3em}
\newgeometry{left=2.5cm,right=2.5cm,top=3cm,bottom=2.5cm}

\renewcommand{\contentsname}{\hspace*{4em}{\smallСодержание}}
\setlength{\cftbeforesecskip}{0.5em}
\renewcommand{\cftsecleader}{\cftdotfill{\cftdotsep}}
\renewcommand{\cftsecaftersnum}{.\hspace{2mm}}

\title{Отчёт по "Управление ИТ-проектами"}
\author{Выполнил студент группы М-РИТ-25-1, Станиславчук Сергей}

\begin{document}

\setstretch{1.5}
\noindent
\begin{minipage}[t]{0.01\textwidth}
\includegraphics{/home/stanik/programmer/++Programmer/5_2/LSTU.jpg}
\end{minipage}%
\begin{minipage}[c]{1.1\linewidth}
    \centering
    {\small
    МИНИСТЕРСТВО НАУКИ И ВЫСШЕГО ОБРАЗОВАНИЯ РФ \\[1em]
    ФЕДЕРАЛЬНОЕ ГОСУДАРСТВЕННОЕ БЮДЖЕТНОЕ \\
    ОБРАЗОВАТЕЛЬНОЕ УЧРЕЖДЕНИЕ ВЫСШЕГО ОБРАЗОВАНИЯ \\
    «ЛИПЕЦКИЙ ГОСУДАРСТВЕННЫЙ ТЕХНИЧЕСКИЙ УНИВЕРСИТЕТ»
    }
\end{minipage}%

\vspace{3em}

\noindent
\hspace{1em}%
\begin{minipage}[t]{0.8\textwidth}
\textbf{Институт} \hspace{2em} компьютерных наук \\
\textbf{Кафедра} \hspace{2em} автоматизированных систем управления
\end{minipage}


\vspace{6em}

\begin{center}
\textbf{Практическая работа №16} \\
по дисциплине \\ <<Управление ИТ-проектами>>
\end{center}

\vspace{6em}

\noindent
\begin{tabular*}{\textwidth}{@{\extracolsep{\fill}} l c r }
    Студент М-РИТ-25-1 & \makebox[4cm]{\hrulefill} & Станиславчук С.М. \\[2em]
Руководитель & \\
К.т.н., доцент & \makebox[4cm]{\hrulefill} & Ведищев В.В. \\
\end{tabular*}

\vfill
\begin{center}
    Липецк, 2026
\end{center}

\thispagestyle{empty}

\newpage
\tableofcontents
\newpage

\section{Управление процессом тестирования программного продукта}

\subsection{Тестовые сценарии}

Тестовые сценарии представляют собой формализованные последовательности действий, предназначенные для проверки корректности функционирования программного продукта при различных условиях эксплуатации. Для системы корпоративного обмена сообщениями тестовые сценарии охватывают как стандартные пользовательские операции, так и альтернативные и исключительные ситуации, связанные с отказами сети, ошибками авторизации и перегрузкой инфраструктуры.

Основной целью разработки тестовых сценариев является проверка соответствия поведения системы функциональным и нефункциональным требованиям. Для рассматриваемой системы тестирование охватывает процедуры аутентификации, обмен текстовыми сообщениями, передачу файлов, голосовую связь и обработку ошибок.

\begin{table}[H]
\centering
\caption{Основные тестовые сценарии системы}
\begin{tabular}{|p{0.8cm}|p{4cm}|p{3.5cm}|p{4cm}|}
\hline
\textbf{№} &
\textbf{Сценарий} &
\textbf{Входные данные} &
\textbf{Ожидаемый результат} \\
\hline

1 &
Авторизация пользователя &
Корректный логин и пароль &
Успешный вход в систему, выдача токена доступа \\
\hline

2 &
Ошибка авторизации &
Некорректный пароль &
Отказ в доступе и отображение сообщения об ошибке \\
\hline

3 &
Отправка текстового сообщения &
Текст сообщения &
Сообщение доставлено адресату и отображается в интерфейсе \\
\hline

4 &
Передача файла &
Файл размером до 100 МБ &
Файл успешно передан и доступен для загрузки \\
\hline

5 &
Голосовая связь &
Подключение двух клиентов &
Устанавливается соединение с допустимой задержкой передачи аудио \\
\hline

6 &
Восстановление соединения &
Временная потеря сети &
Повторное подключение без потери состояния сессии \\
\hline

7 &
Нагрузочное тестирование &
100 одновременных пользователей &
Система сохраняет работоспособность без критических задержек \\
\hline

\end{tabular}
\end{table}

Представленные сценарии позволяют проверить как базовую функциональность, так и устойчивость системы к отказам и нагрузке.

\subsection{Тестовые данные}

Тестовые данные представляют собой наборы входной информации, используемые для проверки поведения системы в различных режимах эксплуатации. Для обеспечения полноты тестирования используются корректные, ошибочные, граничные и нагрузочные данные.

При тестировании авторизации используются корректные и некорректные учётные записи, пустые значения, длинные строки и попытки повторного ввода. Для проверки обмена сообщениями применяются короткие и длинные текстовые сообщения, специальные символы, вложения и бинарные данные.

При тестировании передачи файлов используются документы различных форматов, изображения, архивы и файлы большого размера. Для голосовой связи используются записи аудиопотоков различной продолжительности и искусственно создаваемые сетевые задержки, потери пакетов и джиттер.

\begin{table}[H]
\centering
\caption{Примеры тестовых данных}
\begin{tabular}{|p{3.5cm}|p{4.5cm}|p{5cm}|}
\hline
\textbf{Компонент} &
\textbf{Тестовые данные} &
\textbf{Назначение проверки} \\
\hline

Авторизация &
Корректный и неверный пароль, пустые поля &
Проверка доступа и обработки ошибок \\
\hline

Текстовые сообщения &
Короткий текст, длинный текст, специальные символы &
Проверка корректности передачи сообщений \\
\hline

Передача файлов &
Файлы от 1 КБ до 1 ГБ &
Проверка стабильности передачи и ограничений \\
\hline

Голосовая связь &
Аудиопоток, искусственная задержка, потеря пакетов &
Проверка устойчивости голосового канала \\
\hline

Нагрузка &
100–500 одновременных запросов &
Оценка производительности сервера \\
\hline

\end{tabular}
\end{table}

Использование разнообразных тестовых данных позволяет повысить полноту покрытия тестирования и выявить ошибки, возникающие только в определённых условиях.

\subsection{Тестовое программное обеспечение}

Для тестирования системы используется набор специализированных инструментальных средств, позволяющих автоматизировать запуск тестов, анализировать производительность и проверять устойчивость системы.

В процессе тестирования серверной части используются модульные тесты на языке C++, средства автоматической сборки и тестирования, а также инструменты анализа логов и сетевого взаимодействия. Для контроля качества клиентской части применяются средства проверки интерфейса и пользовательских сценариев.

Для нагрузочного тестирования применяются инструменты генерации сетевого трафика и многопользовательских соединений, позволяющие имитировать одновременную работу большого количества клиентов.

\begin{table}[H]
\centering
\caption{Используемое тестовое программное обеспечение}
\begin{tabular}{|p{4cm}|p{4cm}|p{5cm}|}
\hline
\textbf{Средство} &
\textbf{Назначение} &
\textbf{Область применения} \\
\hline

Git &
Контроль изменений &
Отслеживание исправлений и тестируемых версий \\
\hline

CMake + Make &
Автоматизация сборки &
Проверка успешности компиляции проекта \\
\hline

Google Test &
Модульное тестирование &
Проверка логики серверной части \\
\hline

Wireshark &
Анализ сетевого трафика &
Контроль передачи данных и протоколов \\
\hline

iperf &
Нагрузочное тестирование сети &
Измерение пропускной способности \\
\hline

Valgrind &
Поиск ошибок памяти &
Анализ утечек памяти и ошибок доступа \\
\hline

\end{tabular}
\end{table}

Применение тестового программного обеспечения позволяет автоматизировать процесс контроля качества и уменьшить вероятность пропуска критических ошибок.

\subsection{Тестовое аппаратное обеспечение}

Тестовое аппаратное обеспечение представляет собой совокупность вычислительных и сетевых ресурсов, используемых для проверки функционирования системы в условиях, близких к реальной эксплуатации.

Для тестирования предусматривается отдельный серверный узел, выполняющий функции серверной части системы, а также несколько клиентских рабочих станций, подключённых через локальную сеть. Дополнительно используются маршрутизаторы, коммутаторы и средства имитации сетевых задержек.

\begin{table}[H]
\centering
\caption{Тестовое аппаратное обеспечение}
\begin{tabular}{|p{4cm}|p{3cm}|p{5cm}|}
\hline
\textbf{Оборудование} &
\textbf{Количество} &
\textbf{Назначение} \\
\hline

Сервер &
1 &
Размещение серверной части системы \\
\hline

Клиентские ПК &
3--10 &
Имитация работы пользователей \\
\hline

Коммутатор LAN &
1 &
Организация локального сетевого взаимодействия \\
\hline

Маршрутизатор &
1 &
Тестирование сетевой маршрутизации \\
\hline

Гарнитуры / микрофоны &
2--5 &
Проверка голосовой связи \\
\hline

Накопители данных &
Несколько &
Тестирование передачи файлов \\
\hline

\end{tabular}
\end{table}

Использование специализированного тестового аппаратного обеспечения позволяет оценивать производительность, устойчивость и качество работы системы в условиях, максимально приближенных к корпоративной среде эксплуатации.

\end{document}
